% ==========================================================================
%  SPHERExLens -- presentation specification: every slide, every figure and
%  the complete spoken script. 30 slides; full route 28:10, 15-minute route 18:40.
%  Layout follows the reference "slides and script" document:
%  title page + deck-map table, then one page per slide with
%  "On the slide." / "Figure / diagram." / the figure / "Spoken script."
%  Build:  pdflatex slides_and_script.tex   (twice)
% ==========================================================================
\documentclass[11pt,a4paper]{article}

\usepackage[T1]{fontenc}
\usepackage[utf8]{inputenc}
\usepackage{lmodern}
\usepackage{microtype}
\usepackage[a4paper,margin=2.0cm,top=2.0cm,bottom=2.0cm]{geometry}
\usepackage{amsmath,amssymb}
\usepackage{booktabs,array,longtable,multirow}
\usepackage{graphicx}
\usepackage{xcolor}
\usepackage{enumitem}
\usepackage{titlesec}
\usepackage{fancyhdr}
\usepackage{tikz}
\usetikzlibrary{arrows.meta,positioning,shapes.geometric,calc,fit,backgrounds,patterns,decorations.pathreplacing}
\usepackage{pgfplots}
\pgfplotsset{compat=1.18}
\usepackage{listings}
\usepackage[hidelinks]{hyperref}

% --------------------------------------------------------------- colours ---
\definecolor{ink}{HTML}{14213A}
\definecolor{slate}{HTML}{3F4A5A}
\definecolor{muted}{HTML}{7B8798}
\definecolor{accent}{HTML}{1F6FB2}
\definecolor{accent2}{HTML}{E8F1F8}
\definecolor{good}{HTML}{1E7B4F}
\definecolor{good2}{HTML}{E7F3EC}
\definecolor{bad}{HTML}{B3261E}
\definecolor{bad2}{HTML}{FAECEA}
\definecolor{warn}{HTML}{8A5A00}
\definecolor{warn2}{HTML}{FBF3E2}
\definecolor{rule}{HTML}{C9D0DA}
\definecolor{purple}{HTML}{7C3AED}
\definecolor{purple2}{HTML}{F3E8FF}

% ------------------------------------------------------------ page style ---
\pagestyle{plain}
\setlength{\parindent}{0pt}
\setlength{\parskip}{4pt}

% ---------------------------------------------------------------- macros ---
\newcommand{\arcmin}{\ensuremath{{}^{\prime}}}
\newcommand{\arcsec}{\ensuremath{{}^{\prime\prime}}}
\newcommand{\code}[1]{\texttt{\small\detokenize{#1}}}

% slide header:  Slide 1 --- Title       0:45
% The header is measured first and scaled down to the text width when a long
% title would otherwise run into the running time on the right.
\newsavebox{\slideheadbox}
\newcommand{\slidehead}[3]{%
  \sbox{\slideheadbox}{\large\bfseries Slide #1 --- #2\hspace{1em plus 1fill}#3}%
  \noindent\ifdim\wd\slideheadbox>\linewidth
    \resizebox{\linewidth}{!}{\usebox{\slideheadbox}}%
  \else
    \makebox[\linewidth][s]{\usebox{\slideheadbox}}%
  \fi\par
  \vspace{2pt}\hrule height 0.9pt\vspace{8pt}}

% figure wrapper: sized to fit cleanly without overflowing margins
\newcommand{\figbox}[1]{\begin{center}\resizebox{0.98\linewidth}{!}{#1}\end{center}}


\usepackage[active,tightpage]{preview}
\PreviewEnvironment{tikzpicture}
\begin{document}
\begin{tikzpicture}[font=\footnotesize]
  % NEP Orbit Geometry on Left
  \draw[rule,thick,fill=white,rounded corners=2pt] (0.2,1.2) rectangle (6.2,5.0);
  \node[anchor=north west,font=\scriptsize\bfseries,text=ink] at (0.35,4.85) {SPHEREx Polar Orbit Geometry};

  \draw[slate,thick] (3.2,3.1) circle (1.05);
  \fill[ink] (3.2,3.1) circle (0.12);
  \node[font=\tiny\bfseries,text=white] at (3.2,3.1) {N};
  \node[font=\fontsize{5}{6}\selectfont,text=ink,anchor=north] at (3.2,2.95) {North Ecliptic Pole};

  \foreach \ang in {30, 75, 120, 165, 210, 255, 300, 345}{
    \draw[-{Stealth[length=3pt]},accent,thick] ({3.2+1.05*cos(\ang)},{3.1+1.05*sin(\ang)}) -- ++({-0.22*sin(\ang)},{0.22*cos(\ang)});
  }
  \node[draw=accent,fill=accent2,rounded corners=2pt,align=center,inner sep=2.5pt,text width=5.6cm,font=\tiny] at (3.2,1.6) {
    \textbf{$\sim$20{,}800 Visits to Ecliptic Poles}: Roll shifts band per visit.
  };

  % Lightcurve panels on Right
  \begin{scope}[shift={(6.8,0.35)}]
    % Raw lightcurve (unnormalized)
    \draw[rule,thick,fill=white,rounded corners=2pt] (0,2.6) rectangle (9.2,4.65);
    \node[anchor=north west,font=\tiny\bfseries,text=bad] at (0.2,4.5) {Raw Photometry (Apparent variation driven by LVF roll angle)};
    \draw[bad,thick] plot[smooth] coordinates {
      (0.4,3.2) (1.2,4.0) (2.0,3.0) (2.8,3.9) (3.6,3.1) (4.4,4.1) (5.2,3.2) (6.0,3.9) (6.8,3.1) (7.6,4.0) (8.6,3.2)
    };

    % Spectrum-normalized lightcurve (true variable)
    \draw[rule,thick,fill=white,rounded corners=2pt] (0,0.45) rectangle (9.2,2.4);
    \node[anchor=north west,font=\tiny\bfseries,text=good] at (0.2,2.25) {Spectrum-Normalized: True Variable Isolated (Gaia DR3 confirmed)};
    \draw[good,thick] plot[smooth] coordinates {
      (0.4,1.3) (1.2,1.8) (2.0,1.0) (2.8,1.8) (3.6,1.0) (4.4,1.8) (5.2,1.0) (6.0,1.8) (6.8,1.0) (7.6,1.8) (8.6,1.0)
    };
    \node[draw=good,fill=good2,rounded corners=2pt,inner sep=2pt,font=\tiny,anchor=south east] at (9.0,0.6) {Lomb-Scargle: $P = 1.42$\,d};
  \end{scope}

  \node[draw=rule,thick,fill=white,rounded corners=2pt,align=left,inner sep=3.5pt,text width=15.4cm,font=\tiny] at (8.0,0.1)
    {\textbf{NEP Validation (\code{spherexlens.deep}):} Forced photometry across 78 visits divided by each star's median spectrum with a 3\% systematic floor. Exactly 1 of 6 stars identified as variable --- matching Gaia DR3's independent catalogued variable star with zero false alarms.};
\end{tikzpicture}
\end{document}
